\section{Related Work}\label{sec:related}

\paragraph{Network Motifs in Complex Systems.}
Network motifs were introduced by \citet{milo2002network} as recurring subgraph patterns that appear significantly more often in real networks than in degree-preserving randomized ensembles. The feedforward loop (FFL) emerged as the most functionally characterized motif in biological gene regulatory networks \cite{shenorr2002network}, where it implements sign-sensitive delay and noise filtering \cite{mangan2003structure, alon2007network}. The type-1 coherent FFL, in which the direct and indirect paths have the same sign, acts as a persistence detector, while the type-1 incoherent FFL acts as a pulse generator \cite{alon2007network}. Beyond gene regulation, motif analysis has been applied to technological networks \cite{milo2004superfamilies}, social networks \cite{leskovec2010signed}, ecological food webs \cite{stouffer2007evidence}, and neuronal connectivity in \emph{C.\ elegans} \cite{milo2002network}. \citet{milo2004superfamilies} showed that networks from similar domains share characteristic motif profiles (``superfamilies''), suggesting that motif composition reflects functional requirements. Our work extends this tradition to neural network attribution graphs, testing whether artificial neural circuits exhibit the same structural regularities found in biological and technological systems.

\paragraph{Mechanistic Interpretability and Circuit Discovery.}
The circuits framework \cite{elhage2021mathematical, olsson2022context} provides a conceptual foundation for understanding how transformer networks implement computations through identifiable subnetworks. Key findings include induction heads for in-context learning \cite{olsson2022context}, indirect object identification circuits \cite{wang2023interpretability}, and greater-than circuits \cite{hanna2023how}. Automated circuit discovery methods have scaled beyond manual analysis: ACDC uses activation patching \cite{conmy2023automated}, edge attribution patching identifies minimal circuits \cite{syed2023attribution}, and sparse feature circuits trace computation through sparse autoencoder features \cite{marks2024sparse}. \citet{templeton2024scaling} scaled sparse autoencoders to large language models, enabling attribution graph construction at unprecedented scale. While these methods identify circuits for individual tasks, they do not address whether recurring \emph{structural patterns} exist across tasks. Our motif analysis bridges this gap by identifying cross-task regularities.

\paragraph{Attribution Graphs and the Neuronpedia Platform.}
Sparse autoencoders \cite{bricken2023monosemanticity, cunningham2023sparse} decompose neural network activations into interpretable features, forming the basis for attribution graph construction. Neuronpedia \cite{neuronpedia2024} provides a public repository of attribution graphs that trace causal influence between sparse autoencoder features for specific model behaviors. Each graph is a weighted directed acyclic graph (DAG) where nodes represent features across transformer layers and edges represent signed attribution scores quantifying causal influence. Prior analyses have focused on individual graphs for qualitative circuit interpretation; our work conducts the first systematic quantitative analysis across 200 graphs spanning 8 capability domains, enabling statistical claims about structural universality.

\paragraph{Weighted and Signed Network Motifs.}
\citet{onnela2005intensity} introduced intensity and coherence measures for weighted triangle motifs, extending binary motif analysis to capture quantitative edge information. The intensity of a motif instance is defined as the geometric mean of its edge weights, while coherence measures the homogeneity of weights. Subsequent work extended weighted motif analysis to directed networks \cite{fagiolo2007clustering} and temporal networks \cite{paranjape2017motifs}. We adapt these measures to attribution graphs, adding domain-specific features: sign-coherence (fraction of edges with matching signs), path dominance (ratio of indirect to direct attribution), and weight asymmetry --- all tailored to signed, layered DAGs where edge signs carry semantic meaning.

\paragraph{Graph-Level Statistics for Neural Network Analysis.}
Graph-level statistics (degree distributions, density, clustering coefficients, spectral properties) have been used to compare neural network architectures \cite{you2020design}, characterize training dynamics \cite{frankle2019lottery}, and predict generalization \cite{jiang2019fantastic}. These global descriptors capture important structural properties but may miss local circuit-level patterns. Our unique information decomposition (\S\ref{sec:results_h4}) explicitly disentangles motif-level from graph-level information, showing that they carry complementary structural signals about capability domains.

\paragraph{Higher-Order Motifs and Compositional Structure.}
While 3-node motifs have dominated the literature, several studies have examined higher-order motifs (4-node and beyond) in biological and social networks \cite{milo2004superfamilies}. \citet{benson2016higher} introduced higher-order organization based on network motifs as building blocks for community detection, showing that motif-based clustering reveals structure invisible to traditional methods. In neuroscience, 4-node motifs in cortical connectivity have been linked to functional specialization \cite{sporns2004motifs}. However, whether higher-order motifs in artificial neural networks are compositionally derived from simpler motifs --- specifically, whether universal 4-node patterns are fully embedded within 3-node FFL structures --- has not been previously investigated. Our analysis of 4-node motif FFL-derivativeness (\S\ref{sec:results_h5}) addresses this gap, establishing the compositional hierarchy of motif structure in attribution circuits. Understanding whether complex motifs decompose into simpler building blocks has implications for the parsimony of structural descriptions: if higher-order motifs are fully FFL-derivative, the 3-node FFL provides a sufficient basis for characterizing circuit architecture at multiple scales.

\paragraph{Node Importance and Ablation in Networks.}
Identifying structurally important nodes is a fundamental problem in network science. Traditional centrality measures (degree, betweenness, eigenvector, PageRank) quantify node importance from different perspectives \cite{freeman1978centrality, bonacich1987power}. In the context of neural networks, ablation studies have been used to assess the functional importance of individual neurons \cite{morcos2018importance}, attention heads \cite{voita2019analyzing, michel2019sixteen}, and circuit components \cite{wang2023interpretability}. Our approach differs by defining node importance through motif participation --- specifically, the number of FFL instances containing a node --- and validating this definition through graph-theoretic ablation with multiple matched controls (\S\ref{sec:results_h2}).